% Real-run slide (spec, slide 5): picker × utility dot plot, median of 30 runs per cell, with SPY and equal weight.
% The hedge is defined on the utility slide, so the legend only names it. The paper-trading note is on the takeaways
% slide (user, 2026-10-30). The plot is as wide as the text area allows and its x range is trimmed to the dots, so
% close dots spread out. The legend is stacked in the plot's upper right, where the 33-name rows leave it empty, so the
% plot has the legend's line in height (user, 2026-10-30: mock-up 1, the axis keeps its title). The footer line is
% empty, so the space above it is cut and the plot runs down toward the page number.
% Review round (user, 2026-10-30): an academic title in place of "What you hold beats how you weight it"; each row is
% named a selection rule, as the legend names each dot a weighting rule.
% Row labels in plain words (user, 2026-11-03, mock-up C): each says how many names it holds and, for the 13-stock
% rows, the pool it draws from, so "all 373" no longer reads as holding 373; set in code/src/TalkData.jl. The plot is
% 0.35cm narrower than before so the longest label, "Top 13 by score, 100 most traded", stays inside the margin.
{\renewcommand{\vnfootskip}{0.2em}
\begin{frame}{Stock selection matters more than the weighting rule}
  % a group, not braces: beamer would read a brace group after the title as a frame subtitle
  \begingroup\small Real-data backtest, Jan 2025 – Sep 2026, 373 names. Each row is a selection rule.\par
  \endgroup
  \vspace{7pt}
  \vndotreadrefs{data/realrun_refs.csv}%
  \pgfplotsset{vndot extra/.style={extra description/.code={\vndotlegendstacked}}}%
  \vndotplot[xmin=0.97, xmax=1.72, xstep=0.1, spy=\vndotSPY, ew=\vndotEW, width=9.75cm, height=4.77cm,
    xlabel={terminal wealth ÷ initial wealth, median of 30 runs}]{data/realrun_grid.csv}{data/realrun_bars.csv}\par
  % speaker notes (user, 2026-10-31: algorithm cheat sheets), two pages. Sources: arxiv-paper/sections/
  % method_allocation.tex (selection rules, bandit) and method_mpc.tex (backtest design), verified against
  % code/src/Screens.jl, Bandit.jl, StudyRunner.jl; config/study.toml [screens], [grid].
  \note{\footnotesize
    \vnnotehead{Selection rules (paper's names in parentheses)}
    \vnnoterow{In each sector}{3 names in each of the 11 sectors, 33 in all: the highest γ (γ per sector), the
      bandit's top 3 (bandit per sector; momentum, below), or a uniform random draw.}
    \vnnoterow{Top 13 by score}{The 13 highest γ in the 373-name universe, no sector quota (global γ).}
    \vnnoterow{Diversified}{Walk down the γ ranking; accept a name only if its 252-day growth correlation with every
      accepted name is at most 0.7; fill any shortfall from the ranking (correlation-aware γ).}
    \vnnoterow{100 most traded}{The 100 names with the largest median daily dollar volume over 63 days; the 13
      highest γ among them (liquid then γ) or 13 at random.}
    \vnnoterow{Random 13}{Drawn uniformly from those 100 or from all 373, afresh at every pick.}
    \vnnoterow{My 13 stocks}{AAPL, MSFT, INTC, MU, AMD, NVDA, GS, BAC, WFC, C, F, GM, JNJ, chosen in advance from
      sectors I follow, without 2025 data; the second list adds WMT and PG. Never re-picked (hand-picked).}
    \vnnoterow{Pick dates}{The first day, every 21 trading days and every band exit, with data dated on or before
      the pick. Each pick also refreshes the hedge pool and the short candidates of the 130/30 rule.}}
\end{frame}
\note{\footnotesize
  \vnnotehead{The bandit, and how the runs are read}
  \vnnoterow{Bandit arms}{One ε-greedy bandit per sector, retrained from scratch at each pick; each of the sector's
    15 to 60 names is an arm.}
  \vnnoterow{A pull}{Draws a 21-day window inside the trailing year and earns the name's growth rate over that
    window minus SPY's; $\min(\max(50A, 500), 5000)$ pulls for $A$ arms.}
  \vnnoterow{Exploration}{$\epsilon_k = \max\{0.05,\, \min[1,\, k^{-1/3} (A \ln k)^{1/3}]\}$, which stayed above
    0.51, so every pull explored more than half the time. The rule keeps the 3 arms with the highest mean
    reward.}
  \vnnoterow{In effect}{A sector-neutral ranking by 12-month momentum against SPY.}
  \vnnoterow{Seeds}{30 runs per cell, seeds 5660 to 5689. For the deterministic rules the seed moves only the
    forecast paths, and through them the band exits and the re-plan dates.}
  \vnnoterow{A dot}{The median terminal wealth ratio of a cell's 30 runs, never pooled across rules.}
  \vnnoterow{Benchmarks}{SPY bought and held from the close of 2 Jan 2025; the equal-weight universe bought once
    ($W_0/373$ in each name) and held. No costs.}}
% speaker notes, page 3 (user, 2026-11-03): what hedge on and off mean, for the filled and open dots. Sources:
% code/src/HedgePool.jl (beta_cap, the pool), StudyBacktest.jl, AllocatorsLongShort.jl; config/study.toml [hedge];
% slides/03-utility.tex notes (cap range 1.80 to 0.10 over the window); the 50 of 70 pairs and median gap of
% −0.033 are from data/realrun_grid.csv (hedge true minus false, same row and allocator).
% Synthetic, same pairing (data/syn_*_dots.csv): slow bear (bear stress) higher in 65 of 70, median +0.029;
% crash (burst stress) higher in 34 of 70, median −0.001; typical lower in 59 of 70, median −0.006.
\note{\footnotesize
  \vnnotehead{Hedge on and off (open and filled dots)}
  \vnnoterow{Say}{``With the hedge on, each re-plan puts an upper bound on the portfolio's β, set from the market
    state. The optimizer decides how to get under it. Between re-plans we hold the book, so its β is not
    controlled.''}
  \vnnoterow{Hedge on}{Adds $\beta^{\top} w \le \bar\beta$, $\bar\beta = \max(0,\, 1 - \xi_t)$, and a hedge pool. Hedge off has
    neither; everything else is the same.}
  \vnnoterow{The cap}{1 at $\xi = 0$, looser when bullish ($\xi < 0$), tighter when bearish, 0 at $\xi \ge 1$. Over
    the window it ran from 1.80 down to 0.10.}
  \vnnoterow{Hedge pool}{Up to 5 names with online and down-day β $\le 0$, refreshed at every pick. With
    $\gamma \le 0$ they earn no utility, so they are bought only when a cap needs them, as little as possible.}
  \vnnoterow{Meeting the cap}{The cap is one more constraint; the solver tilts toward lower-β names, holds cash
    (β = 0) or buys pool names (β $\le 0$), whichever costs the least utility.}
  \vnnoterow{Shorts}{Only 130/30 can short, and only names with $\gamma < 0$, held for their own utility term. A
    short of a positive-β name also lowers β, so the solver uses shorts before pool names, up to the 30\% cap.}
  \vnnoterow{On this slide}{Each weighting rule ran both ways. In a mostly rising 2025--26 the hedge usually cost a
    little: lower in 50 of 70 pairs, by a median 0.033 of $W_0$. In the steady decline it helps (higher in 65 of 70); in
    the jump bursts it is a wash.}}
}
